\begin{abstract}
Image-based fashion design with AI techniques has attracted increasing attention in recent years. 
We focus on the reference-based fashion design task, where we aim to combine a reference appearance image and a clothing image to generate a new fashion clothing image. 
% It is a challenging task since there are no reference images available for the newly designed output fashion images. 
Although existing diffusion-based image translation methods have enabled flexible style transfer, it is often difficult to transfer the appearance of the image realistically during reverse diffusion. When the referenced appearance domain greatly differs from the source domain, it often leads to the collapse in the translation.
To tackle this issue, we present a novel diffusion model-based unsupervised structure-aware transfer method, namely \textit{DiffFashion}.
% , to semantically generate new clothes from a source clothing image and a reference appearance image. 
Our method is free of model tuning and structure-preserving and has high flexibility in transferring from images with large domain gaps.
Specifically, based on the optimal transport properties, we keep a shared latent across the clothing image and reference appearance image to bridge the gap between the two domains in the denoising process, and the latent of the reference image is gradually adapted to the clothing domain. 
Simultaneously, the structure is transferred from the source clothing to the output fashion image with mixed guidance, including pre-trained Vision Transformer (ViT) guidance and a foreground mask guidance, to further preserve the structure and appearance semantics from source and reference images. 
% The proposed structure transfer strategy differs from the commonly used diffusion-based translation method where the source image and the style image are in very different latent space. 
% In addition, we employ a mixed guidance in the denoising process, including a foreground mask guidance and pre-trained Vision Transformer (ViT) guidance, to further preserve the structure and appearance semantics from source and reference images.
% To ensure that the denoising process transfers the structure to match the clothing image, we employ a mixed guidance, including a foreground mask guidance and pre-trained Vision Transformer (ViT) guidance. 
% In addition, we enhance the transfer of the structure pattern by mask guidance, utilizing automatically generated semantic masks based on conditioned labels from the diffusion model.
Our experimental results show that the proposed method outperforms state-of-the-art baseline models, generating more realistic images in the fashion design task. Code and demo can be found at \href{https://github.com/Rem105-210/DiffFashion}{https://github.com/Rem105-210/DiffFashion}.
\end{abstract}

%"The utilization of AI techniques in image-based fashion design has been the subject of growing interest in recent years. Our focus is directed towards an innovative fashion design challenge where the objective is to transfer the aesthetic of a reference image onto a clothing image while safeguarding the latter's structural integrity. This poses a formidable task as there is a lack of reference images for the newly designed fashion outputs. Although current diffusion-based image translation or neural style transfer methods allow for adaptable style transfer, it remains a challenge to retain the original image structure in a convincing manner, particularly when the reference image differs significantly from the typical appearance of clothing. To surmount this issue, we introduce a novel unsupervised structure-aware transfer approach based on diffusion modeling, which semantically creates new garments from a given clothing image and reference appearance image. The foreground clothing is separated through the utilization of automatically generated semantic masks, conditioned by labels. These masks are then employed as a guide in the denoising process to preserve the structural information. Additionally, we employ a pre-trained vision Transformer for both appearance and structure guidance. Our experiments demonstrate that our proposed method outperforms existing state-of-the-art models, producing more lifelike images in the fashion design task."